\documentclass[11pt]{article}

\usepackage[utf8]{inputenc}
\usepackage[T1]{fontenc}
\usepackage{lmodern}
\usepackage{amsmath}
\usepackage{amssymb}
\usepackage{amsthm}
\usepackage[hidelinks]{hyperref}
\hypersetup{
  unicode,
  pdftitle={Projected Yukawa Operator and Chiral Polar Factor:
            The Squared Observable HΠ = YΠ†YΠ and the Undetermined Chiral Polar Map UΠ},
  pdfauthor={Jérôme Beau},
}
\usepackage{doi}

\input{glyphtounicode}
\pdfgentounicode=1

\newtheorem{theorem}{Theorem}
\newtheorem{proposition}{Proposition}
\newtheorem{lemma}{Lemma}
\newtheorem{corollary}{Corollary}
\theoremstyle{definition}
\newtheorem{definition}{Definition}
\newtheorem{hypothesis}{Hypothesis}
\theoremstyle{remark}
\newtheorem{remark}{Remark}

\newcommand{\JPi}{J_{\Pi}}
\newcommand{\Jthree}{J_{3}}
\newcommand{\Cgen}{\mathbb{C}^{3}_{\mathrm{gen}}}
\newcommand{\Epi}{E_{\Pi}}
\newcommand{\Sbundle}{\mathcal{S}_{\Pi}}
\newcommand{\Eweak}{E_{\mathrm{weak}}}
\newcommand{\GL}{{\Cgen}^{(L)}}
\newcommand{\GR}{{\Cgen}^{(R)}}
\newcommand{\LY}{L_{Y}}
\newcommand{\Ypi}{Y_{\Pi}}
\newcommand{\Hpi}{H_{\Pi}}
\newcommand{\Upi}{U_{\Pi}}
\newcommand{\lamY}{\lambda_{Y}}
\newcommand{\Sym}{\mathrm{Sym}}
\newcommand{\SU}{\mathrm{SU}}
\newcommand{\UU}{\mathrm{U}}
\newcommand{\diag}{\operatorname{diag}}
\newcommand{\Spec}{\operatorname{Spec}}
\newcommand{\rank}{\operatorname{rank}}

\title{
  Projected Yukawa Operator and Chiral Polar Factor
  \\[1ex]
  \large The Squared Observable $\Hpi = \Ypi^{\dagger}\Ypi$ and the Undetermined Chiral Polar Map $\Upi$
}

\author{J. Beau\\
Independent Researcher, France\\
\texttt{jerome.beau@cosmochrony.org}}

\date{2026}

\begin{document}

  \maketitle

  \begin{abstract}
    This note takes the second step of the fermionic-matter mass-sector frontier.
    The companion line note types the Yukawa line as $\LY = \wedge^{2}(\Eweak)$ under [H-Spin] and [H-Weak] of Q14, and
    reads generation levels on the model operator $\diag(1, \tfrac12 + u, \tfrac12 - u)$ of $\Cgen$.
    Three premises, none supplied by a source, are named:
    [H-Res], the model operator is $\Epi^{2}|_{\Cgen}$;
    [H-Sq], $\Hpi = \Ypi^{\dagger}\Ypi = \lamY^{2}\,\Epi^{2}|_{\Cgen}$;
    [H-Fac], the full coupling factorises through its generation block.
    The results concern the generation block $\Ypi$ and transfer to the full operator only under [H-Fac].
    The polar decomposition $\Ypi = \Upi\,\Hpi^{1/2}$ is linear algebra; under [H-Res] and [H-Sq],
    $\Spec(\Hpi) = \lamY^{2}\{1, \tfrac12 + u, \tfrac12 - u\}$, and the unitary $\Upi$, which carries the chiral
    orientation and the mixing, is invisible to the squared residue.
    Its class $[\Upi]$ for the $\Jthree$-commutants of the carriers is non-trivial exactly when its generator has a
    transverse part, absent on the real cascade: no mixing.
    On the orientation-compatible branch of Q14, an input, the real diagonal split is a Lorentz-chirality
    imbalance, not the $V{-}A$ structure of [H-Weak].
    Q14 excludes a complex metaplectic phase as a source of the external block $R_{\mathrm{mix}}$ within its model;
    the internal block $e_{0} \leftrightarrow e_{\pm}$ is neither excluded by Q14 nor constructed here.
    The static spectrum fixes an order-one level split, not the charged-fermion hierarchy.
  \end{abstract}

  \noindent\textbf{Keywords:}
  projected Yukawa operator; hermitian square; polar decomposition; chiral polar factor; polar class;
  rephasing invariants; Jarlskog phase; generation mixing; generation levels; partial no-go; Schur residue;
  CP-real branch; non-injective projection.

  \newpage

  \tableofcontents

  \newpage

  \section{Position and purpose}
    \label{sec:position}

    The line note~\cite{Beau2026pyl} types the coupling line conditionally.
    It types it as the determinant line $\LY = \wedge^{2}(\Eweak)$ of the distinct rank-two weak factor $\Eweak$
    of~\cite{Beau2026q14}, which carries the abelian weight only because the structure group of $\Eweak$ is $\UU(2)$,
    and it reads generation levels on the model operator $\diag(1, \tfrac12 + u, \tfrac12 - u)$ of $\Cgen$, as levels of
    $\Epi^{2}|_{\Cgen}$ under the premise [H-Res] of Section~\ref{sec:morphism}.
    The factor $\Eweak$ and the Lorentz reading of the spinor factor $\Sbundle$ are two hypotheses
    of~\cite{Beau2026q14}, [H-Weak] and [H-Spin], supplied by no source; under them the left fermions lie in
    $P_{L}\Sbundle \otimes \Eweak \otimes \LY^{k}$ and the right fermions in $P_{R}\Sbundle \otimes \LY^{m}$.
    The mass matrix is a distinct object: the projected Yukawa morphism between these sectors, whose norm, the
    sign of $u$, and the mixing phase remain open.

    This note analyses that morphism, and the analysis begins with a negative result.
    The squared projective residue is a hermitian object on a single Lorentz chirality, while $\Ypi$ ties the two
    chiralities; the chiral connecting datum is therefore not visible to $\Epi^{2}$.
    Concretely, under [H-Res] and [H-Sq] (Hypotheses~\ref{hyp:identification} and~\ref{hyp:square}),
    $\Epi^{2}|_{\Cgen}$ determines the hermitian square $\Hpi = \Ypi^{\dagger}\Ypi$, hence the squared Yukawa levels,
    but leaves $\Ypi$ free up to a unitary chiral polar factor.
    The purpose of the note is to make this precise, to extract exactly what is nonetheless fixed --- the
    spectrum of $\Ypi^{\dagger}\Ypi$ --- and to delimit what is not.
    Two kinds of statements are kept apart.
    The algebraic results use only the generation block $\Ypi : \GL \to \GR$, positive hermitian squares and the
    $\Jthree$ grading of the two generation carriers; the spinor, weak and line factors are spectators of the block.
    The premises that are not algebraic are named: [H-Res] and [H-Sq] for the reading of the model operator as the
    Yukawa square, [H-Fac] (Hypothesis~\ref{hyp:factorisation}) for the transfer from the block to the full coupling,
    and the branch input of the Lorentz chirality (Section~\ref{sec:morphism}).
    Each statement is tagged below by the premises it uses.
    The discipline of the projective-residue reduction is kept throughout: $\Epi = -\Pi_{S} D (1 - P) D
    \Pi_{S}^{*} = -M^{\dagger} M$ is negative semi-definite, with the split $u$ carried by the transported
    chiral defect~\cite{Beau2026prs}, so a semi-definite operator must not be promoted by hand to an
    oriented chiral morphism.

  \section{The Yukawa operator as a chiral morphism}
    \label{sec:morphism}

    Under the hypotheses [H-Spin] and [H-Weak] of~\cite{Beau2026q14}, the left and right fermion sectors of the line
    note~\cite{Beau2026pyl} are
    \[
      \mathcal{F}_{L} = P_{L}\Sbundle \otimes \Eweak \otimes \LY^{k},
      \qquad
      \mathcal{F}_{R} = P_{R}\Sbundle \otimes \LY^{m},
      \qquad
      \LY = \wedge^{2}(\Eweak),
    \]
    each tensored with the generation factor $\Cgen$, on which the gauge group acts trivially (Theorem~5.2
    of~\cite{Beau2026q14}).
    The projected Yukawa coupling is a morphism between these sectors.
    Write $G$ for the structure group of the bundles: the spin group $\mathrm{Spin}(3,1)$ of [H-Spin] acting on the
    spinor factors, and $\UU(2)$ of [H-Weak] acting on $\Eweak$ and, through the determinant, on $\LY$ (colour factors
    omitted).
    The full coupling lies in $\mathrm{Hom}_{G}(\mathcal{F}_{L} \otimes \Cgen, \mathcal{F}_{R} \otimes \Cgen)$.
    Its component on the generation factor, the generation block, is
    \begin{equation}
      \label{eq:yukawa-morphism}
      \Ypi : \GL \longrightarrow \GR ,
    \end{equation}
    from the left generation carrier $\GL$ to the right generation carrier $\GR$.
    Both are copies of $\Cgen = \Sym^{2}(V_{\mathrm{gen}})$, each with its own Cartan generator
    $\Jthree = \diag(0, 1, -1)$ on the weight basis $(e_{0}, e_{+}, e_{-})$.
    The block is the full coupling only under the premise [H-Fac] below.

    The mass matrix is the value of $\Ypi$ on the locked configuration, and the two physically invariant
    hermitian objects are the left and right squares $\Ypi^{\dagger}\Ypi$ on $\GL$ and
    $\Ypi\,\Ypi^{\dagger}$ on $\GR$.
    The squared projective residue $\Epi^{2}$ is a hermitian operator on a single chirality.
    Under [H-Res] and [H-Sq] it is read as the left square up to the norm, and it cannot by itself carry
    the chiral connecting map between the two squares.
    This is the structural origin of the polar ambiguity below.

    The chiralities in this paper are Lorentz chiralities of the spinor factor under [H-Spin].
    They are not the electroweak chirality of $V{-}A$ selection, which is the representation content of [H-Weak] and
    which this note neither uses nor derives.

    In this paper $\JPi$ denotes the internal antilinear parity of the generation copy, an operator on $\Cgen$ built
    from the Born--Infeld datum; it does not use [H-Spin].

    \begin{hypothesis}[{[H-Res]}: the model operator on $\Cgen$ is the restriction of $\Epi^{2}$; not supplied]
      \label{hyp:identification}
      The model operator $\diag(1, \tfrac12 + u, \tfrac12 - u)$ on $\Cgen = \Sym^{2}(V_{\mathrm{gen}})$, with
      $V_{\mathrm{gen}}$ the distinct supplied generation doublet of~\cite{Beau2026q14} and $v = 0$, is the
      restriction to $\Cgen$ of the square $\Epi^{2}$ of the projective residue $\Epi$ of the projected Dirac
      operator, which exists under [H-Spin].
    \end{hypothesis}

    In~\cite{Beau2026q14} (Section~6) these operators are model operators on $\Cgen$, and their reading as
    restrictions of $\Epi^{2}$ is an identification that is not supplied; the Schur reduction~\cite{Beau2026prs}
    leaves it as an open front.
    [H-Res] is a missing identification, not a refutation.
    It concerns the identity of the model operator and nothing else.
    Every statement that reads the model operator as $\Epi^{2}|_{\Cgen}$ is tagged [H-Res].

    \begin{hypothesis}[{[H-Sq]}: the Yukawa square is the residue square; not supplied]
      \label{hyp:square}
      With $\lamY > 0$ a normalisation, the left Yukawa square of the generation block is
      $\Hpi := \Ypi^{\dagger}\Ypi = \lamY^{2}\,\Epi^{2}|_{\Cgen}$.
    \end{hypothesis}

    [H-Sq] is a dictionary between the Yukawa square and the residue square.
    It is a separate premise from [H-Res] and is used together with it: [H-Res] fixes what $\Epi^{2}|_{\Cgen}$ is,
    [H-Sq] equates it with the Yukawa square.
    No source supplies either.
    The level operator itself, $\diag(1, \tfrac12 + u, \tfrac12 - u)$, and the linear algebra built on it do not
    depend on them.

    \begin{hypothesis}[{[H-Fac]}: the full coupling factorises through the generation block; not supplied]
      \label{hyp:factorisation}
      Let $A = \mathcal{F}_{L}$, $B = \mathcal{F}_{R}$, and let $G$ act trivially on $\Cgen$.
      Then $\mathrm{Hom}_{G}(A \otimes \Cgen, B \otimes \Cgen) = \mathrm{Hom}_{G}(A, B) \otimes \mathrm{End}(\Cgen)$.
      [H-Fac] requires that $\mathrm{Hom}_{G}(A, B)$ be spanned by one fixed intertwiner $\Phi$: it is non-zero
      (existence) and one-dimensional (uniqueness up to scale).
      The full coupling is then $\Phi \otimes \Ypi$, with $\Ypi \in \mathrm{Hom}(\GL, \GR)$ the generation block.
    \end{hypothesis}

    Equivalently, the multiplicity space of the pairing of the $G$-isotypic components of $A$ and $B$ is
    one-dimensional: $\sum_{\rho} m_{A}(\rho)\,m_{B}(\rho) = 1$, the sum running over the irreducible $G$-types $\rho$,
    with $m_{A}(\rho)$, $m_{B}(\rho)$ the multiplicities of $\rho$ in $A$ and $B$.
    Existence needs a common irreducible $G$-type of $A$ and $B$.
    As typed, $A$ is built on the $\UU(2)$-doublet $\Eweak \otimes \LY^{k}$ and $B$ on the $\UU(2)$-character
    $\LY^{m}$, which share no irreducible type, and $P_{L}\Sbundle$ and $P_{R}\Sbundle$ are inequivalent modules of the
    spin group.
    So $\mathrm{Hom}_{G}(A, B) = 0$ until a carrier linking the weak doublet to the weak singlet, with the matching
    Lorentz pairing, is supplied.
    Uniqueness is a further condition on the multiplicities of that supplied datum.
    The corpus constructs neither: $\Eweak$ is not constructed (Hypothesis~2.3 of~\cite{Beau2026q14}), nor is the
    linking carrier, nor an intertwiner $\Phi$.
    This is a missing element, not a refutation: nothing proves that the carrier cannot be supplied.
    The results below are stated for the generation block.
    They transfer to the full operator only under [H-Fac]; without it, the polar decomposition, the rephasing class
    and the transverse criterion are statements about the block.

    Branch input.
    Left-admissibility, $P_{R}\Epi P_{R} = 0$, equivalent for a faithful Weyl symbol to $\pi_{LL} = 0$ for the
    eliminated block~\cite{Beau2026prs}, holds on the orientation-compatible branch of Q14 (Proposition~3.8
    of~\cite{Beau2026q14}).
    That branch is an input: it is conditional on [H-Spin] and on the Born--Infeld datum (H1)--(H2)
    of~\cite{Beau2026a34}, and the spinorial Born--Infeld lift does not select it.
    Left-admissibility is not derived here.
    It is a statement about Lorentz chirality, compatible with but not selecting the left-handed weak assignment of
    [H-Weak] (the $V{-}A$ structure), and the two are not identified here (Remark~3.9 of~\cite{Beau2026q14}).
    Every use of it below is tagged ``branch input''.

  \section{The squared Yukawa observable}
    \label{sec:squared}

    \begin{proposition}[Squared Yukawa observable]
      \label{prop:squared-yukawa}
      Let $\Ypi : \GL \to \GR$ be the generation block of the projected Yukawa morphism, and assume [H-Res] and [H-Sq]
      (Hypotheses~\ref{hyp:identification} and~\ref{hyp:square}).
      Then the generation-level observable determined by the Schur-residue sector is the positive hermitian
      operator
      \[
        \Hpi := \Ypi^{\dagger}\Ypi
      \]
      on the left generation carrier.
      On the gauge-singlet triplet $\Cgen$, its dimensionless spectral part is fixed, up to the overall
      Yukawa norm $\lamY$, by
      \[
        \Hpi\big|_{\Cgen} = \lamY^{2}\,\diag\!\left(1,\ \tfrac12 + u,\ \tfrac12 - u\right) ,
      \]
      so that $\Spec\big(\Ypi^{\dagger}\Ypi\big) = \lamY^{2}\,\{1,\ \tfrac12 + u,\ \tfrac12 - u\}$.
      Thus $\Epi^{2}|_{\Cgen}$ determines the squared Yukawa levels, not the Yukawa morphism itself.
    \end{proposition}

    \begin{proof}
      The Schur reduction gives $\Epi = -M^{\dagger} M$ negative semi-definite~\cite{Beau2026prs}, so $\Epi^{2}$ is
      positive semi-definite.
      By [H-Res], the model operator $\diag(1, \tfrac12 + u, \tfrac12 - u)$, the normal form of~\cite{Beau2026prs} with
      $v = 0$ whose even part $(C_{2} - \Jthree^{2})/C_{2} = \diag(1, \tfrac12, \tfrac12)$ is algebraic, is
      $\Epi^{2}|_{\Cgen}$.
      By [H-Sq], $\Hpi = \Ypi^{\dagger}\Ypi = \lamY^{2}\,\Epi^{2}|_{\Cgen}$ with $\lamY$ the undetermined overall norm,
      which gives $\Hpi|_{\Cgen} = \lamY^{2}\,\diag(1, \tfrac12 + u, \tfrac12 - u)$.
      The eigenvalues are the diagonal entries, so the spectrum is as stated, and a hermitian square fixes a
      hermitian operator, not the underlying morphism, which is the content of Section~\ref{sec:polar}.
    \end{proof}

    Premises: [H-Res] and [H-Sq] for the identification, the generation block, and positive hermitian squares;
    the spinor, weak and line factors are spectators of the block, and [H-Fac] is needed to read $\Hpi$ as the left
    square of the full operator.

    The positive root of $\Epi^{2}$ on its domain is $-\Epi$ itself, since $\Epi$ is negative semi-definite.
    [H-Res] concerns $\Epi^{2}$ only: it supplies no restriction of $\Epi$ to $\Cgen$, so $-\Epi|_{\Cgen}$ is not read
    as the positive root $\Hpi^{1/2}/\lamY$ of the model operator.
    No oriented chiral information has been used, in keeping with the reduction discipline.

  \section{Polar decomposition and the undetermined chiral phase}
    \label{sec:polar}

    \begin{proposition}[Chiral polar ambiguity]
      \label{prop:polar-ambiguity}
      The operator $\Hpi = \Ypi^{\dagger}\Ypi$ determines $\Ypi$ only up to a polar factor:
      \[
        \Ypi = \Upi\,\Hpi^{1/2} ,
      \]
      where
      \[
        \Upi : \GL \longrightarrow \GR
      \]
      is a unitary chiral polar map on the support of $\Hpi$.
      Consequently, under [H-Res] and [H-Sq], the squared residue fixes the generation levels but not
      the complex mixing phase, nor the full Yukawa operator.
    \end{proposition}

    \begin{proof}
      For a positive hermitian $\Hpi$ with $\Hpi^{1/2}$ its positive square root, any $\Ypi$ with
      $\Ypi^{\dagger}\Ypi = \Hpi$ satisfies $\Ypi = \Upi\,\Hpi^{1/2}$ with $\Upi := \Ypi\,\Hpi^{-1/2}$ a
      partial isometry, unitary on the support of $\Hpi$, since $\Upi^{\dagger}\Upi = \Hpi^{-1/2}
      \Ypi^{\dagger}\Ypi\,\Hpi^{-1/2} = \Hpi^{-1/2}\Hpi\,\Hpi^{-1/2} = \mathbf{1}$ on that support.
      Conversely, for any such $\Upi$ the product $\Upi\,\Hpi^{1/2}$ has square $\Hpi^{1/2}\Upi^{\dagger}
      \Upi\,\Hpi^{1/2} = \Hpi$, so the solution set is exactly $\{\Upi\,\Hpi^{1/2}\}$.
      The map $\Upi$ acts into the right carrier $\GR$, and $\Hpi$ is invariant under $\Ypi \mapsto V\,\Ypi$
      for every unitary $V$ on $\GR$, $(V\Ypi)^{\dagger}(V\Ypi) = \Ypi^{\dagger} V^{\dagger} V\,\Ypi =
      \Ypi^{\dagger}\Ypi$, so distinct $\Upi$ give distinct $\Ypi$ with identical $\Hpi$.
      Hence $\Hpi$, equivalently (under [H-Res] and [H-Sq]) $\Epi^{2}|_{\Cgen}$, does not
      determine $\Ypi$: the generation-mixing
      content and the relative phases reside in $\Upi$, on which $\Epi^{2}$ has no hold.
    \end{proof}

    Premises: the generation block and positive hermitian squares only; the first assertion is linear algebra and
    uses none of [H-Res], [H-Sq], [H-Fac].
    [H-Res] and [H-Sq] enter only the reading of $\Hpi$ as $\Epi^{2}|_{\Cgen}$, and [H-Fac] the transfer to the full
    operator.

    \begin{remark}[The partial no-go, sharply]
      \label{rem:no-go}
      Proposition~\ref{prop:polar-ambiguity} is a partial no-go: $\Epi^{2}$, read as $\Hpi$ under
      [H-Res] and [H-Sq], does not fix $\Ypi$.
      What survives is exactly the spectral datum of Proposition~\ref{prop:squared-yukawa},
      $\Spec(\Ypi^{\dagger}\Ypi) = \lamY^{2}\{1, \tfrac12 + u, \tfrac12 - u\}$.
      The undetermined factor is not a numerical residue but a unitary chiral map, the carrier of the chiral
      orientation, the relative generation phases, and the mixing.
      The positive scale $\lamY$ is a separate normalisation of the hermitian square $\Hpi$, invisible to the
      dimensionless level operator $\Epi^{2}|_{\Cgen}$, and it is not carried by the unitary $\Upi$.
    \end{remark}

    \begin{remark}[The open datum in $u$ is the orientation branch, not the scalar sign]
      \label{rem:orientation}
      In an oriented $\Jthree$-labelled generation basis $(e_0, e_+, e_-)$ the level operator
      $\diag(1, \tfrac12 + u, \tfrac12 - u)$ already sees the relative sign of $u$.
      The operator $\Hpi$ therefore fixes the sign of $u$ only after such an oriented basis has been chosen
      (under [H-Res] and [H-Sq] for the reading as $\Hpi$).
      Without that chiral-orientation choice the exchange $e_+ \leftrightarrow e_-$ sends $u \mapsto -u$ and
      leaves the unordered squared spectrum unchanged.
      Thus the remaining datum is the physical orientation branch of the split, not an additional scalar
      visible to $\Epi^{2}$ alone.
    \end{remark}

  \section{Diagonal CP-real branch}
    \label{sec:cp-real}

    The reduction~\cite{Beau2026prs} imposes $v = 0$ on the CP-even real sector, so that the off-diagonal mixing
    component $R_{\mathrm{mix}}$ of the model operator vanishes there.
    In the metaplectic step model of~\cite{Beau2026q14} (Section~6, Remark~6.4) no generator reaches $R_{\mathrm{mix}}$,
    whether real or complex.
    In the present language the real branch is $\Upi = \mathbf{1}$.
    There the Yukawa block is, under [H-Res] and [H-Sq], the positive root itself,
    \begin{equation}
      \label{eq:cp-real}
      \Ypi^{\mathrm{CP}} = \Hpi^{1/2} = \lamY\,\diag\!\left(1,\ \sqrt{\tfrac12 + u},\ \sqrt{\tfrac12 - u}\right),
    \end{equation}
    real, diagonal in the generation basis, with no mixing.
    This is a choice of branch, not a derived fact: a non-trivial $\Upi$ --- a generation rotation or a
    relative phase --- leaves $\Hpi$ and the levels untouched while injecting off-diagonal and complex
    structure that $\Epi^{2}$ cannot register.
    The mixing and the CP phase therefore live entirely in $\Upi$, and two routes to them are left open, with
    different statuses.
    The internal block $e_{0} \leftrightarrow e_{\pm}$ is sourced by a non-central $\mathfrak{sl}_{2}$ coefficient with
    $\Im p \neq 0$ or $\Im q \neq 0$ (Proposition~\ref{prop:reality-structure}); Remark~6.4 of~\cite{Beau2026q14}
    concerns the direction $R_{\mathrm{mix}}$, so Q14 does not exclude this route, and none of it is supplied here.
    The external block $e_{+} \leftrightarrow e_{-}$, that is $R_{\mathrm{mix}}$, is excluded as a source within the
    step model of~\cite{Beau2026q14} and needs a generator outside the image of $\mathfrak{sl}_{2}(\mathbb{C})$.

  \section{Polar class of the projected Yukawa operator}
    \label{sec:polar-class}

    The polar factor $\Upi$ of Section~\ref{sec:polar} is not a canonical observable before one quotients by
    the admissible chiral rephasings.
    The admissible rephasings of a carrier are the unitaries that preserve its $\Jthree$ grading, that is, the unitaries
    commuting with $\Jthree = \diag(0, 1, -1)$.
    The weights $0, +1, -1$ are distinct, so these are the diagonal unitaries.
    They are defined by the carrier, a copy of $\Cgen$ with its own grading, and not by $\Ypi$.
    The level operator $\Hpi$ is positive and commutes with $\Jthree$; for $0 < u < \tfrac12$ its spectrum
    $\lamY^{2}\{1, \tfrac12 + u, \tfrac12 - u\}$ (under [H-Res] and [H-Sq]) is distinct, so on the left carrier its
    eigenbasis is the weight basis.
    The right carrier carries the level operator $\Ypi\,\Ypi^{\dagger} = \Upi\,\Hpi\,\Upi^{\dagger}$, which has the same
    spectrum as $\Hpi$ as for any square matrix, and whose eigenbasis is the right weight basis rotated by $\Upi$.
    The class of $\Upi$ measures this misalignment of the right level eigenbasis with the right $\Jthree$ grading.
    The antiunitary exchange of the two carriers of~\cite{Beau2026prs} gives the same spectral reading and is not
    used.

    \begin{proposition}[Polar class]
      \label{prop:polar-class}
      Let $\Hpi = \Ypi^{\dagger}\Ypi$ be positive, commuting with $\Jthree$, with distinct eigenvalues, as
      $\lamY^{2}\,\diag(1, \tfrac12 + u, \tfrac12 - u)$ with $0 < u < \tfrac12$ is under [H-Res] and [H-Sq].
      Let $V_{L} \in \UU(1)^{3}$ act on $\GL$ and $V_{R} \in \UU(1)^{3}$ on $\GR$ as the diagonal rephasings of the
      weight bases, that is the unitaries commuting with $\Jthree$ on each carrier.
      Then $\Ypi \mapsto V_{R}\,\Ypi\,V_{L}^{-1}$, $\Hpi \mapsto V_{L}\,\Hpi\,V_{L}^{-1} = \Hpi$, and
      \[
        \Upi \longmapsto V_{R}\,\Upi\,V_{L}^{-1} ,
      \]
      so the projected Yukawa construction carries not a preferred matrix $\Upi$ but its rephasing class, the
      double quotient
      \[
        [\Upi] \in \UU(1)^{3}_{R} \,\backslash\, \UU(3) \,/\, \UU(1)^{3}_{L} .
      \]
      The class $[I]$ is the set of diagonal unitaries.
      Its rephasing invariants are the moduli $|(\Upi)_{ij}|$ together with the Jarlskog-type phase
      \[
        J_{\Pi} = \Im\!\big( (\Upi)_{11}(\Upi)_{22}\,\overline{(\Upi)_{12}}\,\overline{(\Upi)_{21}} \big) .
      \]
      The common diagonal phase acts trivially on the double action, so the class space has dimension
      $9 - (2\cdot 3 - 1) = 4$, namely three mixing angles and one CP phase, and the differential of the moduli and
      $J_{\Pi}$ has rank four at generic points.
    \end{proposition}

    \begin{proof}
      The commutant of $\Jthree$, a diagonal matrix with the distinct entries $0, 1, -1$, consists of diagonal
      matrices, since $(W\Jthree - \Jthree W)_{ij} = W_{ij}({\Jthree}_{jj} - {\Jthree}_{ii})$.
      The groups $\UU(1)^{3}_{L}$ and $\UU(1)^{3}_{R}$ are therefore fixed by the gradings of the two carriers, with
      no reference to $\Ypi$ or to the eigenbasis of $\Ypi\,\Ypi^{\dagger}$.
      The positive operator $\Hpi$ commutes with $\Jthree$ and has distinct eigenvalues, so it is diagonal in the weight
      basis and commutes with $V_{L}$ and with $\Hpi^{-1/2}$.
      Hence $\Hpi' = V_{L}\,\Ypi^{\dagger}\Ypi\,V_{L}^{-1} = \Hpi$ and
      $\Upi' = V_{R}\,\Ypi\,V_{L}^{-1}\Hpi^{-1/2} = V_{R}\,\Upi\,V_{L}^{-1}$.
      The choice $V_{R} = \Upi$ would give $V_{R}^{-1}\,\Ypi = \Hpi^{1/2}$ and remove $\Upi$ entirely.
      But $\Upi$ commutes with $\Jthree$ only if it is diagonal, so it is not an admissible rephasing.
      Pinning $\GR$ by the eigenbasis of $\Ypi\,\Ypi^{\dagger}$ would be this choice, and the right group is not
      defined that way.
      Under rephasings $(\Upi)_{ij} \mapsto e^{i(a_{i} - b_{j})}(\Upi)_{ij}$, so the moduli are invariant and the
      quartet phase of $J_{\Pi}$ changes by $(a_{1} - b_{1}) + (a_{2} - b_{2}) - (a_{1} - b_{2}) - (a_{2} - b_{1}) = 0$.
      The orbit of $I$ is $\{V_{R}V_{L}^{-1}\}$, the diagonal unitaries.
      A rephasing fixing $\Upi$ satisfies $(\Upi)_{ij}(e^{ia_{i}} - e^{ib_{j}}) = 0$, so for a $\Upi$ with no vanishing
      entry, an open dense set, all $e^{ia_{i}}$ and $e^{ib_{j}}$ coincide: the stabiliser is the common phase.
      Orbits there have dimension $2\cdot 3 - 1 = 5$, and the dimension of $\UU(3)$ being nine, four physical parameters
      survive, three angles and one phase.
      The rank of the differential of the moduli and $J_{\Pi}$ is checked in \texttt{front\_polar\_class\_rephasing.py}.
    \end{proof}

    Premises: the generation block, the $\Jthree$ grading of the two carriers (each a copy of $\Cgen$), and a positive
    $\Hpi$ commuting with $\Jthree$ with distinct spectrum; [H-Res] and [H-Sq] enter only through the reading of that
    $\Hpi$ as $\lamY^{2}\diag(1, \tfrac12 + u, \tfrac12 - u)$, and [H-Fac] through the transfer to the full operator.
    The spinor, weak and line factors are spectators of the block.
    The distinctness of the spectrum makes the left group $\UU(1)^{3}_{L}$ equal to the commutant of $\Hpi$.
    At $u = 0$ the commutant of $\Hpi$ is larger than $\UU(1)^{3}_{L}$.
    The class defined here is then finer than the class under that commutant, so the statement is made for
    $0 < u < \tfrac12$.

    The next obstruction is therefore not the existence of $\Upi$, but the non-triviality of its class: whether
    the generator of the step model produces $[\Upi] \neq [I]$ or only a representative diagonal-rephasing
    equivalent to the identity.

    \begin{proposition}[Polar non-triviality]
      \label{prop:polar-nontriviality}
      Write $\Upi(\gamma) = \exp(\gamma A)$ for the chiral generator $A = \partial_{\gamma}\Upi|_{0}$ induced by
      the complex metaplectic phase $\gamma$ of the metaplectic step model, with $A$ anti-hermitian.
      The diagonal-rephasing orbit through the identity has tangent space the diagonal anti-hermitian
      matrices, so $[\Upi(\gamma)] \neq [I]$ at first order if and only if $A$ has a non-zero off-diagonal
      (transverse) part.
      The rephasing-invariant detectors may be taken as
      \[
        |(\Upi)_{ij}|^{2} = \gamma^{2}|A_{ij}|^{2} + O(\gamma^{3}), \qquad i \neq j ,
      \]
      or, equivalently for a positive branch of $\gamma$, $\partial_{\gamma}^{+}|(\Upi)_{ij}|_{\gamma = 0} =
      |A_{ij}|$, so the off-diagonal moduli detect precisely the transverse component of the metaplectic
      generator; the Jarlskog invariant has leading non-vanishing order $\gamma^{3}$ and is proportional to
      $\Im(A_{12}A_{23}A_{31})$, independent of the diagonal phases.
    \end{proposition}

    \begin{proof}
      A curve $V_{R}(s)\,V_{L}(s)^{-1}$ through the identity in the rephasing group has derivative
      $i(D_{R} - D_{L})$ at $s = 0$, diagonal and anti-hermitian, and every diagonal anti-hermitian is so
      reached; hence the orbit tangent at $I$ is exactly the diagonal anti-hermitian subspace.
      For $\Upi(\gamma) = \exp(\gamma A)$ the off-diagonal entry is $(\Upi)_{ij} = \gamma A_{ij} + O(\gamma^{2})$
      so $|(\Upi)_{ij}|^{2} = \gamma^{2}|A_{ij}|^{2} + O(\gamma^{3})$, and the moduli detect precisely the
      transverse part of $A$.
      The quartet phase vanishes at orders $\gamma^{0}, \gamma^{1}, \gamma^{2}$ and its leading $\gamma^{3}$
      coefficient is proportional to $\Im(A_{12}A_{23}A_{31})$, the genuine three-generation datum, with no
      dependence on the diagonal phases.
    \end{proof}

    Premises: the generation block only.
    The statement holds for every smooth curve $\Upi(\gamma) = \exp(\gamma A)$ of unitaries with anti-hermitian $A$; its
    reading in terms of a metaplectic phase is a modelling step of~\cite{Beau2026q14} (Section~6), not a premise of
    the algebra, and [H-Fac] is needed to read the block class as that of the full operator.

    \begin{corollary}[Real-cascade collapse]
      \label{cor:real-cascade-collapse}
      On the real cascade of the model the off-diagonal mixing channel is absent, $v = 0$, the transported chiral
      defect being the real diagonal split~\cite{Beau2026prs}.
      On the branch input the left-admissibility condition reads $\pi_{LL} = 0$, so the split is carried by the
      imbalance $\pi_{RR}$ of the eliminated block: after restriction to $\Cgen$, under [H-Res],
      $u \propto \pi_{RR} - \pi_{LL}$ up to the Weyl transport, refined by the transported defect
      $\Delta_{\chi} = \pi_{LL} - \tau\,\overline{\pi_{RR}}\,\tau^{-1}$ of~\cite{Beau2026prs}, with $\tau$ the unitary
      intertwiner of the chirality exchange.
      This is a statement on Lorentz chirality under [H-Spin] and the branch input, not the $V{-}A$ structure of
      [H-Weak]; the wording ``spontaneous $V{-}A$ choice'' for the sign of $u$ in~\cite{Beau2026prs} is not adopted.
      In a covariant moving family, with a covariant moving spinorial embedding and a $\JPi$-even Dirac operator, the
      first-order coefficient of the external component $R_{\mathrm{mix}}$ of $\Epi^{2}$ vanishes by antiunitary
      parity, and so does every odd-order coefficient; mixing at even orders is not excluded, and the internal entries
      $e_{0} \leftrightarrow e_{\pm}$ are not constrained by that parity argument~\cite{Beau2026ebj}.
      On the real cascade data the generator $A_{\Pi}$ vanishes (Remark~\ref{rem:two-transverse-blocks},
      Proposition~\ref{prop:reality-structure}), so $A$ is diagonal, $\Upi(\gamma)$ stays diagonal, and
      $[\Upi] = [I]$, with $|(\Upi)_{ij}| = \delta_{ij}$ and $J_{\Pi} = 0$.
      No physical mixing is produced at this stratum.
      A non-trivial class $[\Upi] \neq [I]$, hence CKM/PMNS-type mixing, requires a transverse metaplectic
      component.
      In the present real cascade of the model this transverse component is absent.
      The internal block is not excluded by~\cite{Beau2026q14} and is not constructed here; its source would need a
      complex non-central $\mathfrak{sl}_{2}$ coefficient (Remark~\ref{rem:two-transverse-blocks}), which is
      conditional on a phase and an ordering that are not prescribed~\cite{Beau2026bim}.
      The external block is excluded as a source within the step model of~\cite{Beau2026q14}.
      The resulting mixing matrix is physical only as a relative polar class between two fermionic sectors
      sharing a generation carrier.
    \end{corollary}

    Premises: the generation block and the algebra of Proposition~\ref{prop:polar-nontriviality}, together with the
    real cascade data of~\cite{Beau2026prs}; the branch input for the reading of the diagonal split as a
    Lorentz-chirality imbalance, and [H-Res] where the real data are stated as data on $\Epi^{2}|_{\Cgen}$.
    The statement of~\cite{Beau2026ebj} is used under its own hypotheses, and [H-Fac] is needed to read the block class
    as that of the full operator.

    \begin{remark}[Two transverse blocks]
      \label{rem:two-transverse-blocks}
      The transverse obstruction has two distinct components.
      On the generation triplet $\Cgen = \langle e_{0}, e_{+}, e_{-}\rangle$ the polar tangent transverse to the
      diagonal-rephasing orbit decomposes into the internal block $e_{0} \leftrightarrow e_{\pm}$ and the
      external block $e_{+} \leftrightarrow e_{-}$, the latter denoted $R_{\mathrm{mix}}$.
      For a single metaplectic $\mathfrak{sl}_{2}$ cascade lifted to $\Sym^{2}$, the correct unitary chiral
      generator is the anti-hermitian part of the $\JPi$-odd component,
      \[
        A_{\Pi} = \Pi_{\mathrm{antiherm}}\,\Pi_{\JPi\text{-odd}}\big(\dot{\Upi}\big) .
      \]
      On real cascade data this generator vanishes, recovering Corollary~\ref{cor:real-cascade-collapse}.
      The internal block is sourced by the $\Sym^{2}(\mathfrak{sl}_{2})$ lift if the $\mathfrak{sl}_{2}$ coefficients
      have a genuinely complex non-central part.
      Q14 does not exclude this (its Remark~6.4 concerns $R_{\mathrm{mix}}$), and the corpus does not construct such a
      part: it is a condition, not a result.
      By contrast the external block $e_{+} \leftrightarrow e_{-}$ is never sourced by the
      $\Sym^{2}(\mathfrak{sl}_{2})$ lift alone; it lies outside the single-cascade metaplectic image and would
      require an additional tower datum, and a complex metaplectic phase is excluded as a source of it within the
      metaplectic step model of~\cite{Beau2026q14} (Section~6).
      The chiral-frontier normalisation channel and the polar-mixing channel therefore separate: the oriented
      area $N_{A}$ belongs to the real Cartan / $\Jthree$ channel~\cite{Beau2026aar} and can normalise the
      diagonal split without predicting a CKM/PMNS polar class.
      A non-trivial polar class in the present carrier first requires a complex non-central metaplectic
      component, and the external $R_{\mathrm{mix}}$ block requires still more: a generator beyond the single
      $\mathfrak{sl}_{2}$ metaplectic cascade.
    \end{remark}

    \begin{proposition}[Reality structure of the polar generator]
      \label{prop:reality-structure}
      Write the non-central metaplectic step as $M = p\,E + q\,F + r\,H$ in the real $\mathfrak{sl}_{2}$ basis,
      and let $\sigma(M) = \overline{M}$ be the entrywise conjugation whose fixed real form is
      $\mathfrak{sl}_{2}(\mathbb{R}) = \mathrm{span}_{\mathbb{R}}\{E, F, H\}$, so $\sigma(M) = M$ iff
      $p, q, r \in \mathbb{R}$.
      The chiral polar generator $A_{\Pi} = \Pi_{\mathrm{antiherm}}\Pi_{\JPi\text{-odd}}(\dot{\Upi})$ depends only
      on the $\sigma$-anti-real part $\tfrac12(M - \sigma M) = i\,(\Im p\,E + \Im q\,F + \Im r\,H)$, and splits as
      the polar class does.
      Its transverse (off-diagonal) part is non-zero iff $\Im p \neq 0$ or $\Im q \neq 0$, while its diagonal part
      is the pure $\Jthree$ rephasing $i\,\Im r\,\diag(0, 2, -2)$, which lies in the $\UU(1)^{3}$ double-coset
      quotient defining $[\Upi]$ and is therefore class-trivial.
      Hence
      \[
        [\Upi] = [I] \quad\Longleftrightarrow\quad \Im p = \Im q = 0 ,
      \]
      independently of $\Im r$.
      In particular a complex representation matrix obtained from a complex spinorial frame $W$ leaves the class
      unchanged, $A_{\Pi} \mapsto W A_{\Pi} W^{\dagger}$, so the mixing is controlled by the reality of the
      \emph{coefficients} $p, q$, not by that of the matrix entries.
    \end{proposition}

    Premises: the generation block and the $\Sym^{2}(\mathfrak{sl}_{2})$ lift on $V_{\mathrm{gen}}$ of the metaplectic
    step model; the $\JPi$-odd projection uses the internal parity of the generation copy, not [H-Spin].
    [H-Fac] is needed to read the block statement as one on the full operator.

    \begin{corollary}[Vanishing mixing generator at the present stratum]
      \label{cor:negative-closure-present-stratum}
      At the present full-tower stratum the data of~\cite{Beau2026prs, Beau2026aar, Beau2026cho} give a non-central
      step with real coefficients.
      The cascade step is $g = \exp(tE)\exp(sF)$ with $t, s \in \mathbb{R}$, its per-step $\JPi$-odd $\Jthree$
      component being the real oriented area $\alpha = ts$~\cite{Beau2026prs, Beau2026aar}; and the inherited lift
      carries only the central cyclotomic phase $\zeta_{q}^{\Delta A_{c}}$, with no Fourier or Weyl generator and
      hence no metaplectic Gauss phase~\cite{Beau2026cho}, a central datum that is projectively trivial for the
      polar factor.
      Thus $\Im p = \Im q = 0$ at this stratum and, by Proposition~\ref{prop:reality-structure}, $A_{\Pi} = 0$ and
      $[\Upi] = [I]$.
      The diagonal split stands as stratum data: $N_{A} \neq 0$~\cite{Beau2026cho} and $u \neq 0$~\cite{Beau2026prs},
      read as a coefficient of $\Epi^{2}|_{\Cgen}$ under [H-Res].
      The stratum data therefore give a diagonal generation split but no mixing.
      A non-trivial internal block requires a genuinely complex non-central metaplectic phase, a Weil generator
      absent from the present pipeline and available only in a new tower stratum, gated by the same full-tower
      type-rigidity~\cite{Beau2026aog} as the capacity normalisation.
    \end{corollary}

    Premises: the generation block and Proposition~\ref{prop:reality-structure}, together with the real cascade
    and central-phase data of~\cite{Beau2026prs, Beau2026aar, Beau2026cho}, and [H-Res] for the reading of $u$;
    neither the line $\LY$ nor the Lorentz or weak typing enters, and [H-Fac] is needed to read the block statement as
    one on the full operator.

    \begin{remark}[Two delimitations]
      \label{rem:precautions}
      First, the spin-stratum type-rigidity of~\cite{Beau2026aog} removes the orientation caveat
      $[H\text{-orient}]$ within the present spin stratum and excludes a spin-Galois $\sqrt{5}$ factor orthogonal to the
      central phase, but it does not bear on the non-triviality of $[\Upi]$: chiral orientation and polar class are
      distinct data.
      Second, the generation split $u \neq 0$~\cite{Beau2026prs, Beau2026bim} is a diagonal-channel datum and
      produces no mixing; it contributes nothing to the off-diagonal moduli or to $J_{\Pi}$.
      Mixing is the transverse polar datum, not the diagonal split.
    \end{remark}

  \section{Status of results}
    \label{sec:status}

    \paragraph{Derived (this note), for the generation block.}
      The Yukawa block is determined by its hermitian square $\Hpi = \Ypi^{\dagger}\Ypi$ only up to a unitary chiral
      polar factor, $\Ypi = \Upi\,\Hpi^{1/2}$ (Proposition~\ref{prop:polar-ambiguity}), without any of [H-Res], [H-Sq],
      [H-Fac].
      The polar factor is moreover not canonical but a rephasing class for the $\Jthree$-commutants of the two carriers,
      the double quotient $[\Upi] \in \UU(1)^{3}_{R} \backslash \UU(3) / \UU(1)^{3}_{L}$, with invariants the moduli
      $|(\Upi)_{ij}|$ and the Jarlskog phase $J_{\Pi}$, four physical parameters in all
      (Proposition~\ref{prop:polar-class}).
      Its non-triviality is controlled by the transverse component of the metaplectic generator: $[\Upi] \neq
      [I]$ iff that generator has a non-zero off-diagonal component (Proposition~\ref{prop:polar-nontriviality}).
      These propositions are algebraic consequences: they use only the generation block, positive hermitian squares and
      the $\Jthree$ gradings, and hold whatever the typing of the spinor, weak and line factors, with no use of the
      line $\LY$.
      The typing of the fermion sectors, $\LY = \wedge^{2}(\Eweak)$ with the left and right sectors above, is the
      conditional typing of the line note~\cite{Beau2026pyl} under [H-Spin] and [H-Weak] and is carried here only as
      the setting of the spectator factors.

    \paragraph{Conditional.}
      Under [H-Res] and [H-Sq], supplied by no source~\cite{Beau2026q14}, the generation-level observable of the
      Schur-residue sector is the positive hermitian square $\Hpi = \Ypi^{\dagger}\Ypi$, with
      $\Spec(\Hpi)|_{\Cgen} = \lamY^{2}\{1, \tfrac12 + u, \tfrac12 - u\}$
      (Proposition~\ref{prop:squared-yukawa}).
      The transfer of every block statement to the full coupling holds only under [H-Fac], whose existence and
      uniqueness conditions are not met by any constructed object of the corpus (Section~\ref{sec:morphism}).
      On the real cascade of the model the off-diagonal channel is absent, $v = 0$, so $[\Upi] = [I]$ and no
      physical mixing is produced at this stratum (Corollary~\ref{cor:real-cascade-collapse}); the diagonal split is
      read as the Lorentz-chirality imbalance on the branch input, which is an input and not derived, and not as
      $V{-}A$ selection.
      The transverse datum further splits into two blocks (Remark~\ref{rem:two-transverse-blocks}): the internal
      block $e_{0} \leftrightarrow e_{\pm}$ is sourced by a genuinely complex non-central $\mathfrak{sl}_{2}$
      coefficient, a condition that Q14 does not exclude and the corpus does not construct, while the external block
      $e_{+} \leftrightarrow e_{-}$ is not in the image of the single-cascade $\Sym^{2}(\mathfrak{sl}_{2})$ lift and
      is excluded as a source within the step model of~\cite{Beau2026q14}.
      Consequently $N_{A}$ may fix the diagonal oriented-area normalisation without fixing the polar mixing
      class.
      At the present full-tower stratum these non-central coefficients are real by the stratum data --- the cascade
      step is real~\cite{Beau2026prs, Beau2026aar} and the inherited lift carries only the central
      phase~\cite{Beau2026cho} --- so $\Im p = \Im q = 0$, $A_{\Pi} = 0$ and $[\Upi] = [I]$: the stratum data give
      the diagonal split but not the mixing (Proposition~\ref{prop:reality-structure},
      Corollary~\ref{cor:negative-closure-present-stratum}).

    \paragraph{Open.}
      The positive Yukawa norm $\lamY$ and a non-trivial polar class $[\Upi] \neq [I]$ are not fixed here, and the
      premises [H-Res], [H-Sq], [H-Fac] and the construction of $\Eweak$ and of the linking carrier are open.
      The norm $\lamY$ scales $\Hpi$, whereas a non-trivial internal block of $[\Upi]$ requires a genuinely complex
      non-central metaplectic phase, absent from the present pipeline and available only through a new tower
      stratum~\cite{Beau2026aog}, and the external block requires a generator outside the step model.
      It is observable as a mixing matrix only relative to a second fermionic sector.
      Hence the absolute masses, the generation mixing, and the CP structure remain downstream open data, in
      agreement with~\cite{Beau2026fmnote}.

    \paragraph{Arithmetic gate for the diagonal split.}
      The diagonal split parameter \(u\) and the angular--radial datum \(\varepsilon\) are not
      independent invariants.
      They are the same functional of the level ratio
      \[
        R_\Pi
        =
        \frac{\frac12+u}{\frac12-u}.
      \]
      Indeed, the inverse parametrisation
      \[
        \varepsilon(R)
        =
        \frac{R-1}{2(R+1)}
      \]
      gives \(\varepsilon(R_\Pi)=u\) identically.
      This identity is only a functional identification.
      The representation \(\Sym^{2}(\mathbb{C}^{2})\) does not select the value of \(u\):
      the diagonal split
      \[
        \diag\left(1,\frac12+u,\frac12-u\right)
      \]
      commutes with \(\Jthree\) for every \(u\).
      Thus the polar spectrum of \(\Hpi\) (under [H-Res] and [H-Sq]) carries the vertical diagonal split as a free
      carrier-preserving datum, but does not by itself carry the horizontal ADE arithmetic
      obstruction.

      The value \(u=\varepsilon=1/10\) is therefore equivalent to the dictionary gate
      \[
        R_\Pi = R_{\mathrm{ADE}}=\frac32,
      \]
      where \(R_{\mathrm{ADE}}\) is the order-\(5\) Cayley-spectral ratio~\cite{Beau2026aog}.
      This gate is compatible with the closed angular datum, but it is not produced by the
      carrier-preserving recursion.
      In particular, the recursion may fix the vertical parity type and the field
      \(\mathbb{Q}(\sqrt5)\) in which the obstruction is visible, but it does not generate the
      horizontal Galois transport \(2\leftrightarrow 2'\) selecting the order-\(5\) value.
      Consequently the notation \(u\sim_{\mathrm{dict}}\varepsilon\) records the current
      status: same functional, distinct carriers, and value-level equality only after the
      ADE dictionary identification.

    \paragraph{Direct level assignment no-go.}
      Assume the arithmetic dictionary gate \(R_\Pi=R_{\mathrm{ADE}}=3/2\).
      Then \(u=1/10\), and, under [H-Res] and [H-Sq], the squared Yukawa levels fixed by \(\Hpi\) are
      \[
        \lamY^{2}\left\{1,\frac35,\frac25\right\}.
      \]
      Their dimensionless squared ratios are
      \[
        \frac52:\frac32:1,
      \]
      and the corresponding root-level ratios are
      \[
        1:\sqrt{\frac35}:\sqrt{\frac25}.
      \]
      These ratios are of order one.
      They therefore cannot be identified with the observed charged-fermion mass hierarchy by a direct generation
      assignment: the smallest observed inter-generation step already exceeds the whole order-one spread by two orders
      of magnitude.
      The squared projective residue fixes an internal level split, not the physical mass hierarchy.

    \paragraph{Separation of level split and hierarchy.}
      The diagonal datum \(u\) should consequently be read, under [H-Res] and [H-Sq], as a generation-level weight
      inside the positive hermitian square \(\Hpi=\Ypi^{\dagger}\Ypi\), not as a direct mass-ratio predictor.
      The observed hierarchy must be carried by generation-dependent stabilisation depths in the projective
      amplification law, or by a later operator stratum not analysed here.
      It is not carried by a sector-dependent normalisation of \(\lamY\): the Yukawa norm \(\lamY\) remains a single
      downstream scale throughout this note.
      At this stratum, the dictionary gate can fix the order-one split
      \[
        \left\{1,\frac35,\frac25\right\},
      \]
      but it does not fix the charged-fermion hierarchy.
      The charged-lepton mass target is therefore not fixed by the projected Yukawa line itself; it is recorded
      separately as the unit-dispersion Koide target~\cite{Beau2026kud}.

    \paragraph{Hierarchy not carried by the cascade rate.}
      The single-level projective amplification factor is itself of order one: the Kesten--McKay coefficients and the
      ADE representation dimensions give inter-level separations of only a fraction of a log-decade, with spectral
      ratios in \([3/2, 2]\)~\cite{Beau2026a5}.
      Stacked with the order-one split \(\{1, 3/5, 2/5\}\) it therefore stays of order one and cannot reach the
      observed charged-fermion hierarchy.
      The hierarchy is therefore not carried by the order-one split \(u\), nor by the single-level cascade factor
      alone.
      It is carried, if at all, by generation-dependent cascade depths \(n_g\) in the projective amplification law
      \(\Lambda_{\mathrm{proj}}(n)\) in the exponential growth regime~\cite{Beau2026a5}.
      The exponent fixes the amplification rate; the generation stabilisation depths remain the next structural datum.

    \paragraph{Two lemmas: separation and level-crossing cancellation.}
      We record the two structural statements that freeze the present front, both verified by exact symbolic
      computation with no fit (see Section~\ref{sec:repro}).
      Write the mass-channel factorisation in the form
      \begin{equation}
        \label{eq:mass-factorisation}
        m_g \;=\; \lamY\, w_g\, S_g\, A(n_g),
        \qquad A(n_g) = \exp(\beta^{\ast} n_g),
      \end{equation}
      where \(g\in\{1,2,3\}\) is the generation label, \(w_g\) is the Schur-residue weight of the squared
      residue (under [H-Res]), \(S_g\) is the stratigraphic order-one factor, and \(A(n_g)\) is the depth amplification
      generated along the projective cascade at the exponent \(\beta^{\ast}=1/(\delta+\tfrac12)\)~\cite{Beau2026a20}.

      \begin{lemma}[Separation of the generation label and the Schur-residue weight]
        \label{lem:separation-label-weight}
        On the common rank-three carrier \(\Cgen \simeq \mathrm{Sym}^{2}(V_\rho)\) the generation label and the
        Schur-residue weight are distinct, multiplicatively acting factors.
        The label \(g\in\{1,2,3\}\) is fixed by the three admissible stratigraphic levels of the binary
        icosahedral group~\cite{Beau2026a4}.
        The weights \(w_g\in\{1, \tfrac35, \tfrac25\}\) are, under [H-Res] and [H-Sq], the \(J_3\)-triplet of
        \(\Hpi\) at the dictionary gate and are not identified with the stratigraphic level ratios.
      \end{lemma}
      \begin{proof}
        The two factors depend on disjoint data.
        The weight \(w_g\) is a function of the single vertical datum \(u\) alone: the diagonal
        \(\diag(1,\tfrac12+u,\tfrac12-u)\) commutes with \(J_3\) for every \(u\), so the squared-residue split is
        free under \(\mathrm{Sym}^{2}\).
        The stratigraphic factor \(S_g\) is the Kesten--McKay and representation datum \(c_i/\dim\rho_i\) of the
        \(2I\) Cayley spectrum~\cite{Beau2026a5}.
        Their dimensionless ratio-sets differ: \(w_g\) gives \(\{1, \tfrac32, \tfrac52\}\) while the stratigraphic
        levels \(\{20,24,30\}\) give \(\{1, \tfrac65, \tfrac32\}\).
        Hence the two readings cannot be collapsed into a single observable.
      \end{proof}

      \begin{lemma}[Level-crossing depths cannot carry the hierarchy]
        \label{lem:level-crossing-cancellation}
        In the exponential trajectory-branching regime \(\Lambda_{\mathrm{proj}}(n)\simeq\lambda_2(n) =
        L_0\exp(\beta^{\ast} n)\), suppose the generation depths are the level-crossing depths
        \(n_{\mathrm{proj}}(\lambda) = (1/\beta^{\ast})\log(\lambda/L_0)\) at which each stratigraphic level
        \(\lambda_g\) becomes projectable.
        Then the depth amplification of~\eqref{eq:mass-factorisation} satisfies
        \begin{equation}
          \label{eq:level-crossing-cancellation}
          \frac{A(n_i)}{A(n_j)} \;=\; \exp\!\bigl(\beta^{\ast}(n_i-n_j)\bigr) \;=\; \frac{\lambda_i}{\lambda_j},
        \end{equation}
        independently of \(\beta^{\ast}\).
        On the admissible levels \(\{20,24,30\}\) the amplified ratios are therefore bounded by \(3/2\), and the
        level-crossing depths cannot carry the charged-fermion hierarchy.
      \end{lemma}
      \begin{proof}
        Substituting \(n_{\mathrm{proj}}(\lambda) = (1/\beta^{\ast})\log(\lambda/L_0)\) into
        \(A(n)=\exp(\beta^{\ast} n)\) gives \(A(n_{\mathrm{proj}}(\lambda)) = \lambda/L_0\); the exponent cancels
        and the ratio reduces to \(\lambda_i/\lambda_j\).
        For the three admissible levels this is at most \(30/20=3/2\), an order-one quantity.
      \end{proof}

      \begin{corollary}[Intrinsic branching depths]
        \label{cor:intrinsic-branching-depths}
        A charged-fermion hierarchy in~\eqref{eq:mass-factorisation} requires the depths \(n_g\) to be intrinsic
        branching-stabilisation data of the exponential trajectory, decoupled from the stratigraphic level ratios.
        The required gaps \(\Delta n \sim \log(\text{ratio})/\beta^{\ast}\) (of order \(64\) for the \(\tau\!:\!e\)
        step at \(\beta^{\ast}\simeq 0.127\)) are not reachable from the static data
        \(\{\lambda_g,\dim\rho_g,\beta^{\ast}\}\), whose logarithms are of order one.
      \end{corollary}

      \begin{remark}[Scope]
        \label{rem:depths-scope}
        Lemma~\ref{lem:level-crossing-cancellation} does not derive \(n_g\); it proves where \(n_g\) cannot come
        from.
        The hierarchy is thereby reduced exactly to the depth-triplet problem for \(\lambda_2(n)\) along the
        admissible cascade, with the depths isolated as the carrier.
        Deriving an intrinsic branching-stabilisation law from the growth law \(N(\lambda;n)\) of~\cite{Beau2026a5}
        is a separate, later front.
      \end{remark}

    \paragraph{Linear-growth depths do not supply the generation stabilisation scale.}
      We record the negative resolution of the depth front, established by exact computation on the corpus
      growth law with no fit and no mass input (Section~\ref{sec:repro}).
      The required inter-generation depth gap of the exponential law,
      \(\Delta n_{\mathrm{req}} = \log(m_\tau/m_e)/\beta^{\ast} \simeq 64\) at \(\beta^{\ast}\simeq 0.127\),
      is bracketed on both sides by the two natural depths of the linear growth law
      \(N(\lambda;n)\sim F_{\mathrm{KM}}(\lambda)\,n\)~\cite{Beau2026a5}, and neither side reaches it.
      First, the level-crossing depth gives the lower end:
      the amplified ratio collapses to \(\lambda_i/\lambda_j\le 3/2\)
      (Lemma~\ref{lem:level-crossing-cancellation}), a gap of order \(3\).
      Second, the cumulative-capacity cell-entry depth, read on the projected-capacity profiles
      \(\sigma_{\mathrm{pair}}(n),\sigma_c(n)\) of the admissibility cascade, is intrinsic and stable in the
      cascade parameter \(q\) but stays \(\le 10\) and does not split by generation:
      the three blocks of a co-admissible triplet share one capacity profile, hence one cell depth.
      Third, the full saturation rank \(n_{\mathrm{sat},g} = k\dim\rho_g/c_g(p)\), with
      \(c_g(p)=F_{\mathrm{KM}}(\lambda_g)\)~\cite{Beau2026a5}, gives the upper end but fails on stability and
      order:
      the gap \(\Delta n_{\mathrm{sat}}(p)\) diverges with the cascade degree \(p\)
      (from \(83\) at \(p=5\) to \(354\) at \(p=53\)) through the band-edge vanishing \(c_1\to 0\), so
      \(\Delta n_{\mathrm{req}}\) is attained only at a tuned sub-saturation fraction and a single \(p\);
      and the symmetry pin \(c_2=\tfrac12\) fixes \(n_{\mathrm{sat},2}=2k\dim\rho_2\) while
      \(n_{\mathrm{sat},3}=k\dim\rho_3/c_3(p)\) crosses it, so the \(n_2,n_3\) generation order swaps with
      \(p\).
      The structural lock is uniform:
      every depth extracted from the linear law --- crossing, saturation, or any intermediate fraction --- is
      a function of the same coefficients \(c_g(p)\), and therefore inherits both the band-edge divergence
      \(c_1\to 0\) and the symmetry pin \(c_2=\tfrac12\); this holds for the inter-block partition and for the
      stratigraphic inter-level partition \(\{20,24,30\}\) alike.
      A \(q\)-stable, correctly ordered stabilisation depth thus cannot be a depth of
      \(N(\lambda;n)\sim F_{\mathrm{KM}}(\lambda)\,n\):
      the positive law, if it exists, must be a decaying capacity-controlled per-level profile
      \(\sigma_{\lambda_g}(n)\) transported onto the cascade, a distinct object not extracted from the linear
      count.
      This is an honest reduction, not a derivation:
      the depth triplet remains the open carrier (Remark~\ref{rem:depths-scope}), now with the linear-law
      route closed.

    \paragraph{Quantised resolution-density frontier.}
      A distinct modelling route is to quantise the per-level resolution density itself, rather than the
      linear count \(N(\lambda_g;n)\).
      In this formulation the transported profile \(\sigma_{\lambda_g}(n)\) is treated as a Weil--Schur
      resolution-density observable on \(\Cgen\), and the quantum cell is required to be independent of the
      capacity coefficient \(c_g(p)\).
      This removes the band-edge denominator obstruction responsible for
      \[
        n_g=\frac{\Delta I_g}{c_g(p)} ,
      \]
      since the resulting depth gap is flat in \(p\) rather than driven by \(c_1(p)\to 0\).
      This gain does not close the generation-scale problem.
      The quantisation acts on the normalised \(F_{\mathrm{KM}}\)-level, so the forced \(O(1)\) quanta fall in
      the crossing regime rather than in the required stabilisation window.
      The window \(\Delta n\in[50,80]\) is reached only by reintroducing the raw Cayley valency \(|S|=24\), or
      by using a quantum ratio of the same order as the target hierarchy itself.
      The resolution-density route therefore converts the previous obstruction into a sharper question, which
      an exact audit of the transport then settles:
      the \(\mathrm{Sym}^2\) Schur-residue transport is scale-free (homogeneous of degree one), carrying only
      order-one residue ratios and never an absolute magnitude, while \(|S|\) enters solely as the Laplacian
      normalisation \(L=I-A/|S|\) that defines \(F_{\mathrm{KM}}\).
      Re-inserting \(|S|\) thus undoes the normalisation defining the density and quantises the un-normalised
      count, the route already locked by \(n_g=\Delta I_g/c_g(p)\); the valency is not forced, so this short
      static lever closes.

    \paragraph{Decision point.}
      The remaining mass-hierarchy question has been reduced to whether the stabilisation depths \(n_g\) are
      intrinsic dynamical data of the projective cascade, rather than functions of the static ADE,
      Schur-residue, or Kesten--McKay level data.
      All static routes tested so far produce only order-one ratios, or reintroduce the target scale through an
      unforced quantum choice.
      Thus the mass-sector programme reaches a sharp bifurcation:
      either a genuinely intrinsic depth-selection law is derived, or the present stratum must be recorded as
      fixing generation structure and order-one level splitting, but not the charged-fermion hierarchy.

    \paragraph{Dynamic stabilisation status.}
      The remaining dynamic reading of the projective mass cascade does not alter the closure
      of the static sector.
      If the three generation depths are read as parallel thresholds on the same admissibility
      flow, with threshold fractions fixed by the \(J_3=\mathrm{diag}(0,1,-1)\) carrier on
      \(\mathbb{C}^3_{\mathrm{gen}}\simeq\mathrm{Sym}^2(\mathbb{C}^2)\), the resulting depth
      gap is independent of the Kesten--McKay band-edge coefficient.
      In particular, it does not inherit the \(c_1\to 0\) edge divergence or the \(c_2=1/2\)
      symmetry pin that affect all \(c_g\)-routed static depths.
      This is a structural improvement, but not a hierarchical one: the forced gap remains
      on the order-one level-crossing scale.

      The sequential Born--Infeld cycle reading gives the only possible amplification of a
      single intrinsic dynamic number.
      However, every forced dynamic range supplied by the admissibility data remains too small:
      the normalised Heisenberg band, the Alon--Boppana adjacency band, and the generation-level
      span all give cumulative depths far below the charged-fermion hierarchy scale.
      A range large enough to close the hierarchy would have to be introduced either through
      the cascade amplification itself, which makes the cycle count tautological, or through an
      unpinned ceiling-to-floor scale such as \(\beta_{\mathrm{BI}}/A_{\min}\).
      Both routes reintroduce a free normalisation under another name.

      The renewed-floor multi-cycle escape is therefore also closed at this stratum.
      A transfer law
      \[
        \mathrm{floor}_{g+1}=\mathcal{R}(\mathrm{ceiling}_g)
      \]
      would be admissible only if \(\mathcal{R}\) were already forced by the
      Born--Infeld, ADE, Heisenberg, or Schur structure.
      The available forced renewals give only order-one to order-ten cumulative depths; the
      only renewals reaching the charged-fermion scale are circular.
      Thus the projective cascade fixes the three-generation structure, the order-one projected
      splitting, and a band-edge-independent order-one depth gap, but it does not derive the
      charged-fermion hierarchy.
      The hierarchy is not carried by any forced static or first-order dynamic quantity of the
      cascade.

    \paragraph{Premises of the level and depth statements.}
      The statements of this section that read the levels $\lamY^{2}\{1, \tfrac35, \tfrac25\}$ as squared Yukawa
      levels, and the Schur-residue weights $w_g$ of Lemma~\ref{lem:separation-label-weight} as the
      $\Jthree$-triplet of $\Hpi$, use [H-Res] and [H-Sq]; the arithmetic of the level ratios does not.
      Lemma~\ref{lem:level-crossing-cancellation} and Corollary~\ref{cor:intrinsic-branching-depths} concern the
      depth amplification of the exponential law and the static spectral data $\{\lambda_g, \dim\rho_g, \beta^{\ast}\}$;
      they use neither the Yukawa morphism, nor the line $\LY$, nor the Lorentz or weak typing, nor
      [H-Res], [H-Sq], [H-Fac], nor the branch input.

    \paragraph{Summary.}
      Front~3b yields the squared Yukawa levels under [H-Res] and [H-Sq], not a mixing result, and Front~3c
      characterises the chiral polar factor of the generation block as a rephasing class whose non-triviality is the
      transverse metaplectic datum.
      Under [H-Res] and [H-Sq], the squared residue $\Epi^{2}|_{\Cgen}$ fixes $\Hpi = \Ypi^{\dagger}\Ypi$; it does not
      fix $\Ypi$ unless the class $[\Upi]$ is fixed, and that class is the identity on the real cascade of the model.

  \section{Reproducibility}
    \label{sec:repro}

    All operator identities are verified by exact symbolic computation, with no floating-point sampling; the one
    numerical computation, a rank computation in 60-digit arithmetic, is stated where it is used.
    The squared-observable results of Sections~\ref{sec:squared}--\ref{sec:cp-real} ---
    the level substitution and spectrum of Proposition~\ref{prop:squared-yukawa} (the spectrum of $\Ypi^{\dagger}\Ypi$
    for explicit non-trivial unitaries $\Upi$), the polar non-uniqueness of
    Proposition~\ref{prop:polar-ambiguity} (for a generic positive $\Hpi$ and an arbitrary unitary, both a
    generation rotation and a diagonal phase leave $\Ypi^{\dagger}\Ypi$ invariant while changing $\Ypi$),
    the invariance of $\Hpi$ under right-carrier unitaries, and the diagonal no-mixing CP-real branch ---
    are checked in \texttt{simulation/fermionic-matter/front3b\_yukawa\_operator.py}.
    The rephasing class of Proposition~\ref{prop:polar-class} ---
    the $\Jthree$ commutant, the class transformation, the invariance of the moduli and of $J_{\Pi}$ ---
    is checked in exact arithmetic over $\mathbb{Q}(i)$, on a generic and on degenerate $\Hpi$, in
    \texttt{code/front\_polar\_class\_rephasing.py}, whose output is committed beside it.
    The same script computes the rank (four) of the differential of the moduli and $J_{\Pi}$ at random points of
    $\UU(3)$ in 60-digit arithmetic with a fixed seed.
    The remaining polar-class results of Section~\ref{sec:polar-class} ---
    the orbit tangent at the identity, the modulus and Jarlskog
    detectors of the transverse part (Proposition~\ref{prop:polar-nontriviality}), and the diagonal
    real-cascade collapse (Corollary~\ref{cor:real-cascade-collapse}) ---
    are checked in \texttt{front3c\_polar\_class\_audit.py} and
    \texttt{front3c\_polar\_class\_nontriviality.py}.
    The two-block decomposition of Remark~\ref{rem:two-transverse-blocks} ---
    the vanishing of the chiral generator $A_{\Pi}$ on real cascade data, the sourcing of the internal block
    $e_{0} \leftrightarrow e_{\pm}$ by a complex non-central coefficient, the identical vanishing of the
    external block $R_{\mathrm{mix}}$ in the $\Sym^{2}(\mathfrak{sl}_{2})$ lift, and the independence of the
    mixing datum from the real oriented area $N_{A}$ ---
    is checked in \texttt{front3d\_transverse\_route.py}.
    The reality-structure results of Section~\ref{sec:polar-class} ---
    the dependence of $A_{\Pi}$ on the $\sigma$-anti-real part alone, the equivalence
    $[\Upi] = [I] \Leftrightarrow \Im p = \Im q = 0$, the class-triviality of the imaginary $\Jthree$ rephasing,
    the projective triviality of the central scalar, and the spinorial-frame covariance distinguishing matrix
    complexity from coefficient complexity (Proposition~\ref{prop:reality-structure},
    Corollary~\ref{cor:negative-closure-present-stratum}) ---
    are checked in \texttt{front3e\_reality\_structure.py}.
    The separation and level-crossing results of Section~\ref{sec:status} ---
    the distinctness of the Schur-residue weight and the stratigraphic factor on the common rank-three carrier
    (Lemma~\ref{lem:separation-label-weight}), the \(\beta^{\ast}\)-independent cancellation
    \(A(n_i)/A(n_j)=\lambda_i/\lambda_j\) of the level-crossing depths
    (Lemma~\ref{lem:level-crossing-cancellation}), and the unreachability of the required depth gaps from the
    static data (Corollary~\ref{cor:intrinsic-branching-depths}) ---
    are checked in \texttt{front\_depths\_ng\_reconciliation.py}.
    The bracketing of the required depth gap and the closure of the linear-growth route of
    Section~\ref{sec:status} ---
    the over-amplification and band-edge divergence of the saturation rank in the exponential reading, the
    intrinsic but \(\le 10\) and non-splitting cumulative-capacity cell depth on the admissibility-cascade
    profiles, and the \(q\)-instability of both the inter-level gap and the \(n_2,n_3\) order ---
    are checked in \texttt{front\_saturation\_rank\_amplification.py},
    \texttt{front\_capacity\_cell\_depth.py}, and \texttt{front\_interlevel\_cell\_depth.py}.
    The hard-exit valency test closing the quantised resolution-density route is checked in
    \texttt{front\_valency\_forced\_hardexit.py}.
    It verifies that the \(\mathrm{Sym}^2\) Schur-residue transport is scale-free, that the
    \(F_{\mathrm{KM}}\) density is defined on the normalised Laplacian \(L=I-A/|S|\), and that the
    factor \(|S|\) enters only as the count normalisation.
    Consequently re-inserting \(|S|\) quantises the un-normalised count rather than the
    resolution density, returning to the already closed depth law
    \(n_g=\Delta I_g/c_g(p)\).
    The dynamic stabilisation status of Section~\ref{sec:status} ---
    the band-edge independence of the parallel-threshold depth gap and its order-one magnitude, the
    insufficiency of every forced sequential Born--Infeld dynamic range, and the closure of the
    renewed-floor multi-cycle escape under the requirement that the transfer law be forced by the
    Born--Infeld, ADE, Heisenberg, or Schur structure ---
    is checked in \texttt{front\_intrinsic\_stabilisation\_depths.py} and
    \texttt{front\_renewed\_floor\_multicycle.py}.

  \newpage

  \bibliographystyle{plain}
  \bibliography{cosmochrony-bibliography}

\end{document}
